\section{Agent Framework Design and Implementation}
\label{sec:architecture}

PLC Workbench uses compiler and execution feedback to revise generated PLC code. A PLC function block can retain state across scans, so code that compiles can still produce the wrong output on a later scan. We check each submitted version at the interface, compilation, behavior, and bounded verification stages. Figure~\ref{fig:agent-loop} shows where feedback returns to the author.

\begin{figure}[htbp]
  \centering
  \setlength{\unitlength}{1pt}
  \begin{picture}(434,182)
    \put(6,130){\framebox(128,36){\shortstack{Frozen task\\Spec and interface}}}
    \put(134,148){\vector(1,0){12}}
    \put(146,130){\framebox(132,36){\shortstack{Author framework\\Self agent or Codex CLI}}}
    \put(278,148){\vector(1,0){12}}
    \put(290,130){\framebox(128,36){\shortstack{Candidate version\\ST or SCL}}}
    \put(354,130){\vector(0,-1){15}}
    \put(6,78){\framebox(128,36){\shortstack{Bounded LTL check\\PLCverif and CBMC}}}
    \put(146,78){\framebox(132,36){\shortstack{Behavior test\\Frozen test oracle}}}
    \put(290,78){\framebox(128,36){\shortstack{Interface and\\vendor compilation}}}
    \put(290,96){\vector(-1,0){12}}
    \put(146,96){\vector(-1,0){12}}
    \put(70,78){\vector(0,-1){34}}
    \put(6,8){\framebox(128,36){\shortstack{Run receipt\\Metrics and evidence}}}
    \put(164,8){\framebox(248,36){\shortstack{Repairable diagnostic from any check\\Shared Trial enforces the repair budget}}}
    \put(164,26){\vector(-1,0){30}}
    \put(412,26){\line(1,0){16}}
    \put(428,26){\line(0,1){150}}
    \put(428,176){\line(-1,0){216}}
    \put(212,176){\vector(0,-1){10}}
  \end{picture}
  \caption{Both author frameworks submit candidates to the same Trial controller. Each check records evidence. A repairable diagnostic starts another version within the budget.}
  \label{fig:agent-loop}
\end{figure}

\subsection{Authoring and task context}

The Web workbench accepts the task bank and supports single-task generation, checking existing code, and batch runs. Each trial snapshots the task and its verification configuration. The author receives the requirements and fixed interface for the chosen target. The evaluator retains the expected outputs and applies the same judgment rules to both frameworks. In the self agent, the next repair prompt includes the current source and relevant diagnostics. It compresses the tool results into a compact diagnostic payload while keeping the full evidence on disk. Codex CLI receives an isolated workspace with a task file and a checker. Its requests pass through a loopback relay that uses the same configured model service and meters calls through the same Trial. The service uses a configurable Chat Completions base URL, model ID, and API key. The key remains on the host.

The separate interactive console also supports plan mode, local file attachments, and conversation compaction. When the context budget requires compaction, the console asks the model to summarize older exchanges and preserves as much recent original history as the available window allows. Batch trials manage task context through the author packet and diagnostic feedback described above.

\subsection{Platform adapters and repair}

The CODESYS adapter submits ST to the vendor compiler and uses a simulation project to replay scan sequences. The TIA adapter submits SCL through Openness and uses a separate PLCSIM Advanced runtime to execute the behavior tests. Both platforms use the frozen task oracle. A successful functional evaluation can trigger the shared PLCverif and CBMC checks. For supported SCL constructs, the adapter prepares the source model and records the assumptions used by the verifier.

Trial saves every distinct submitted source version, identifies duplicates by hash, and enforces the same limits for both frameworks: six source versions, 24 model requests including transport retries, and 1,800 seconds per task. The agent can repair candidate errors such as a compile rejection or a reproducible behavior failure. It stops on success, exhausted budget, no progress, or an environment error. An unconfirmed source-model counterexample is reported for review. The workbench also supports independent adversarial review for CODESYS and the optional local ST backend. The reviewer proposes scan sequences from the requirements and candidate. Replay records bind to the candidate hash. Reviewer expectations remain proposals for human review, and this optional stage runs outside the batch comparison.

\subsection{Metrics displayed in the workbench}

\begin{description}
  \item[Compilation.] First submission accepted and final candidate accepted are separate counts. Divide each by the planned task count for the group. Report tasks that never reached the compiler alongside these rates.
  \item[Functional testing.] Report passed scenarios divided by all declared scenarios in the group, and show how many ran. A scenario is a sequence of scans. Also count tasks with a successful final workflow outcome (receipt status \texttt{finite\_pass}) over all planned tasks. Preserve separate compile, behavior, and formal statuses so a formal interruption can be distinguished from a behavior failure.
  \item[Bounded formal verification.] Report properties verified within the recorded bound over declared properties. Show violations, unsupported properties, timeouts, and unexecuted checks separately. Report the scan bound and input assumptions alongside the ratio.
  \item[Resource use.] Count actual model requests including retries. Report provider-returned tokens with coverage and missing usage marked. Record elapsed seconds per task and distinguish summed task durations from elapsed batch wall time under concurrency.
\end{description}

The batch results and per-task receipts retain these denominators and evidence. Each platform keeps its own run configuration. The archived TIA experiment pairs self and Codex with \texttt{deepseek-v4.1-flash-1017} and \texttt{qwen3.8-flash}, producing four groups of 32 tasks. Its run receipt records the model settings, task snapshots, and runtime source hashes. The experiment used local source changes, and twelve recorded source hashes differ in content from the current main branch after accounting for line endings. Reproducing the recorded run requires the corresponding runtime source files.
